\documentclass[11pt]{article}

\usepackage[utf8]{inputenc}
\usepackage[T1]{fontenc}
\usepackage{lmodern}
\usepackage{geometry}
\usepackage{amsmath,amssymb,amsfonts,amsthm,mathtools}
\usepackage{enumitem}
\usepackage{microtype}
\usepackage{hyperref}

\geometry{margin=1in}
\hypersetup{colorlinks=true, linkcolor=black, urlcolor=blue, citecolor=black}
\setlist[enumerate]{topsep=0.35em,itemsep=0.3em}
\setlist[itemize]{topsep=0.3em,itemsep=0.25em}

\newcommand{\Aether}{\AE{}ther}
\newcommand{\Aflow}{\AE{}ther-flow}
\newcommand{\TheoryProgram}{\Aether{} / \Aflow{} framework}
\newcommand{\Sub}{\mathcal{A}}
\newcommand{\PhiFlow}{\Phi_{\lambda}}
\newcommand{\Order}{\mathcal{S}}
\newcommand{\Front}{\mathcal{F}}
\newcommand{\Resp}{\mathcal{R}_{\Gamma}}
\newcommand{\Cycle}{\mathcal{C}_{\Gamma}}
\newcommand{\Map}{\mathcal{L}_{\mathrm{loc}}}
\newcommand{\Data}{\mathcal{D}_{\mathrm{loc}}}

\newtheorem{refutationfinding}{Refutation finding}
\newtheorem{blocker}{Promotion blocker}
\newtheorem{repairburden}{Repair burden}
\newtheorem{localverdict}{Local verdict}

\title{\textbf{Localization Repair Refutation}\\[0.35em]
\large Source-Determination Test for the Repaired Localization Packet}
\author{Refuter AgentJob AJ-RT-20260608-006-001}
\date{June 8, 2026}

\begin{document}

\maketitle

\begin{abstract}
This draft attacks the source-repair packet for the localization candidate. The
result is a local negative result, not a global closure of the research
program. The repaired packet improves the earlier candidate by refusing direct
target null-cone, radar-coordinate, and target-proper-time definitions.
However, it still does not determine its key objects from the registered Gate
0 source reservoir. The first-response front depends on an unstated source
order measure and adjacency topology; the coincidence map depends on an
unstated response-mode algebra; and the clock scale assigns local units without
a source transport law. Promotion remains blocked.
\end{abstract}

\section{Scope and Refutation Standard}

This artifact is a draft/control output for task \texttt{RT-20260608-006}. It
does not edit canonical ontology TeX, does not issue a Gate Chair verdict, and
does not reject the broader \TheoryProgram{}. It tests the repaired data
\[
\Data^{+}
=
(\Sub,\PhiFlow,\Gamma,\Order_{\Gamma},\Front^{+},
\kappa^{+}_{\Gamma},\alpha^{+}_{\Gamma},\Map^{+},\mathcal{I}^{+},
\mathcal{T}^{+})
\]
against the source-determination standard stated by the prior audit.

The standard is stricter than verbal anti-smuggling discipline. A repaired
definition passes only if its construction is invariantly specified from
source-side objects before benchmark comparison. A definition fails locally if
it needs an unregistered metric topology, target clock normalization, target
observer protocol, or an adjustable free choice that can later be tuned to the
benchmark.

\section{Front-Generation Findings}

\begin{refutationfinding}[First response uses more than order incidence]
The repaired front defines
\[
\delta_{a\to b}
=
\inf\{\Delta>0:\mathsf{Resp}(R(b\mid a;\lambda_0,\Delta),\epsilon)\}.
\]
This expression is not determined by a bare ordered motion
\(\Phi_{\lambda}\). It requires a quantitative interval \(\Delta\), an
infimum operation, and enough order topology or measure structure to compare
nearby values of \(\Delta\). The Gate 0 reservoir allows an ordering
parameter, but it does not yet supply an invariant source measure on order
intervals.
\end{refutationfinding}

The problem is not notation. The repair simultaneously treats monotone
reparameterization of the source order as redundancy and uses numerical
first-response intervals as defining data. Under a general monotone
reparameterization, the numerical value and limiting behavior of \(\Delta\)
are not invariant unless a transformation law or order-measure primitive has
already been fixed.

\begin{blocker}[Order-measure blocker]
The repaired front is not promotion-ready until the project supplies either a
source-invariant order measure or a purely order-theoretic replacement for
\(\delta_{a\to b}\) that does not depend on a tunable numerical parameter.
\end{blocker}

\begin{refutationfinding}[Boundary minimality imports an unstated adjacency]
The repaired packet defines \(\Front^{+}\) using boundary minimality in a
``source adjacency order.'' The registered Gate 0 formalization supplies
\(\Sub\), \(\Phi_{\lambda}\), an observer-accessible slice placeholder, and an
S-time placeholder. It does not yet supply a source adjacency relation or
topology strong enough to define boundary points.
\end{refutationfinding}

If adjacency is taken as primitive, the repair has added a new source object
without formalizing its invariances or dimensions. If adjacency is derived from
first response, the definition is circular: the front is generated using a
boundary relation that already depends on the response structure the front is
supposed to define.

\begin{blocker}[Adjacency blocker]
The candidate must state whether source adjacency is primitive, derived, or
operational. Each option requires an explicit non-target definition before the
front can be treated as source-generated.
\end{blocker}

\begin{refutationfinding}[The response threshold is a tuning channel]
The response predicate \(\mathsf{Resp}(R,\epsilon)\) and threshold
\(\epsilon\) are not fixed by the source packet. Without a source law for
admissible perturbation classes \(P_a\), response equivalence, and threshold
selection, different choices of \(\epsilon\) can select different effective
fronts. That freedom can be used later to fit the benchmark cone.
\end{refutationfinding}

\section{Coincidence-Response Findings}

\begin{refutationfinding}[Three detector modes are asserted rather than derived]
The repaired packet assigns observer-chain modes
\[
Q_{\Gamma}=(q^0_{\Gamma},q^1_{\Gamma},q^2_{\Gamma},q^3_{\Gamma})
\]
and requires the last three modes to be functionally independent on a local
source-accessible domain. This introduces a \(1+3\) split at the level of the
response protocol. The split is not automatically a target import, but it is
not yet derived from \((\Sub,\PhiFlow,\Order)\).
\end{refutationfinding}

The phrase ``functional independence'' also requires a mathematical structure:
a differentiable response manifold, an algebra of response functionals, a rank
condition, or an equivalent source criterion. The repair states the need for
independence but does not define the structure in which independence is to be
tested.

\begin{blocker}[Response-mode blocker]
Promotion requires a source-only mode-selection law explaining why admissible
observer chains carry exactly three independent non-order response labels and
how independence is checked without Euclidean axes, target simultaneity, or a
preferred target frame.
\end{blocker}

\section{Clock-Scale Findings}

\begin{refutationfinding}[Local cycle units do not define a transported scale]
The repaired clock scale is
\[
\alpha^{+}_{\Gamma}
=
\frac{u_{\Gamma}}{\Delta\Order_{\Gamma}(\Cycle)}.
\]
This assigns a local unit to one stable cycle on one chain. It does not yet
define comparison of units between chains, transport of clock scale through
the source domain, or invariance of the denominator under allowed monotone
reparameterizations of source order.
\end{refutationfinding}

The repair avoids direct target proper-time import, which is a real
improvement. But avoiding import is weaker than determining scale. Unless
\(u_{\Gamma}\), cycle equivalence, bounded drift, and cross-chain comparison
are fixed by source rules, \(\alpha^{+}_{\Gamma}\) remains an adjustable
calibration variable.

\begin{blocker}[Clock-transport blocker]
The candidate needs a source law for cycle equivalence and clock-scale
transport. A local stable cycle can define a counter, but it does not by itself
define the conformal scale of an observer-accessible metric candidate.
\end{blocker}

\section{Benchmark-Separation Finding}

\begin{refutationfinding}[Procedure does not enforce non-tuning]
The repair states that source definitions must be fixed before exact-closure
benchmark comparison. That is necessary. It is not sufficient. A later
constructor could still choose the response threshold, adjacency relation,
mode basis, and clock unit using knowledge of the benchmark, then record the
choice before the formal comparison step.
\end{refutationfinding}

The needed safeguard is a source-side admissibility criterion that makes the
definitions determinate without reference to the target. The present repair
contains an order of operations, not yet a non-tuning theorem.

\section{Local Verdict}

\begin{localverdict}[Promotion blocked by source underdetermination]
The repaired localization packet is not refuted as an explicit target-metric
import. It is refuted as a promotion candidate because its central structures
remain underdetermined by the registered source reservoir. The failure points
are:
\begin{enumerate}
    \item no invariant source-order measure for first-response intervals;
    \item no formal source adjacency or boundary topology;
    \item no source law for response thresholds and perturbation classes;
    \item no source algebra for three independent detector modes; and
    \item no cross-chain clock-scale transport law.
\end{enumerate}
\end{localverdict}

This is a local negative result. It blocks Gate Chair review of the repaired
packet, but it also identifies a precise repair path.

\section{Repair Burdens for the Director}

\begin{repairburden}[Source-basis repair burden]
Before another localization candidate can improve on this result, the project
should formalize a source-basis packet containing:
\begin{enumerate}
    \item an invariant source order measure or an order-only first-response
    replacement;
    \item a source adjacency or topology definition with relabeling
    invariance;
    \item an admissibility law for perturbation classes and response
    thresholds;
    \item a response-mode algebra or rank criterion that explains the local
    \(1+3\) detector split; and
    \item a clock-cycle equivalence and transport law for comparing scale
    across observer chains.
\end{enumerate}
\end{repairburden}

The logical next Director decision is therefore not Gate Chair review. A
bounded Ontology Formalizer job is the most conservative next role if the
Director wants to repair the source basis. A controlled pause is also coherent
if the Director concludes that adding these structures would exceed the
authorized source reservoir.

\section*{References}

Ricciardi, A. S. (2026, June 8). \emph{Localization source-repair packet:
Source-only front, coincidence, and clock-scale definitions} [Unpublished
manuscript]. The Aether-Flow Interpretation of Relativity Research Project,
\texttt{research\_control/tasks/RT-20260608-005/artifacts/04\_LOCALIZATION\_SOURCE\_REPAIR\_PACKET.tex}.

Ricciardi, A. S. (2026, June 8). \emph{Localization smuggling audit:
Target-import review of the Gate 0 localization candidate packet}
[Unpublished manuscript]. The Aether-Flow Interpretation of Relativity
Research Project,
\texttt{research\_control/tasks/RT-20260608-004/artifacts/03\_LOCALIZATION\_SMUGGLING\_AUDIT.tex}.

Ricciardi, A. S. (2026, June 8). \emph{Localization candidate packet: A
source-side observer-localization map for Gate 0 review} [Unpublished
manuscript]. The Aether-Flow Interpretation of Relativity Research Project,
\texttt{research\_control/tasks/RT-20260608-003/artifacts/02\_LOCALIZATION\_CANDIDATE\_PACKET.tex}.

Ricciardi, A. S. (2026, June 8). \emph{Gate 0 ontology formalization: Source
primitives, forbidden imports, and localization burden} [Unpublished
manuscript]. The Aether-Flow Interpretation of Relativity Research Project,
\texttt{research\_control/tasks/RT-20260608-002/artifacts/01\_GATE0\_ONTOLOGY\_FORMALIZATION.tex}.

\end{document}
